\section{Introduction}
Trend estimation is a basic task in time series analysis, but trend smoothness
is rarely known in advance. Standard time-series and forecasting references
emphasize that trend extraction, prediction, and structural interpretation are
related but distinct goals \citep{chatfield2003analysis,brockwell2016introduction,
hamilton1994timeseries,makridakis1998forecasting}. This motivates procedures
that choose smoothness from data while preserving a strict temporal evaluation
protocol. This paper compares the proposed train-validation smoothness
selection approach with a time-weighted validation variant and several
classical or simple benchmarks.

The comparison is implemented as a reproducible subproject that uses
\texttt{trend\_estimation} as a normal Python library dependency
\citep{trendestimation}. The purpose is not to claim universal superiority from
one ticker, but to provide a concrete empirical case study and a template for
broader comparisons. The paper is deliberately explicit about terminology:
Guerrero smoothness is used only for finite-difference penalized smoothers,
whereas realized roughness is reported as a diagnostic for every fitted trend.
